\Def Многоугольник - часть плоскости, ограниченная замкнутой ломанной. В памяти представляется как массив точек в порядке их обхода.

\Def Многоугольник - простой, если в нем отсутствуют самопересечения.

\Def Многоугольник выпуклый, если \begin{enumerate}
    \item любая его диагональ лежит внутри.
    \item любой отрезок с концами в его внутренности лежит внутри целиком.
    \item все его точки лежат по одну сторону от любой прямой, проходящей через сторону.
\end{enumerate}

\Def Многоугольник выпуклый, если векторное произведение соседних
сторон сохраняет знак на протяжении всего обхода.

\begin{lstlisting}
bool IsConvex(Point[] points, int n) {
  double init = 
    Cross(points[1] - points[0], points[2] - points[1]);
  for (int i = 1; i < n; ++i) {
    Point A = points[i];
    Point B = points[(i + 1) % n];
    Point C = points[(i + 1) % n];
      if (Cross(AB, AC) * init < 0) { return false; }
  }
  return true;
}
\end{lstlisting}
\vspace{\baselineskip}
\Task Дан многоугольник(необязательно выпуклый). Найти его площадь.

\textbf{Решение 1(\textit{метод треугольников}):} \newline
Рассмотрим произвольную точку плоскости $A$. Например, одну из
вершин многоугольника. Тогда площадь многоугольника $P_1 ... P_n$
определяется как сумма \textit{ориентированных} площадей треугольников $AP_iP_{i+1}$.

То есть, $S = |\displaystyle\sum_{i=0}^{n-2} S_i|$, где $S_i = \frac{1}{2}\Vec{AP_i} \times \Vec{AP_{i+1}}$\newline

\textbf{Решение 2(\textit{метод трапеций}):} \newline
Площадь многоугольника равна модулю суммы \textit{ориентированных}
площадей трапеций, порождаемых сторонами многоугольника. То есть
$S_i = \displaystyle\frac{(P_{i+1}.x - P_i.x)(P_{i+1}.y+P_i.y)}{2}$
\pagebreak